% -----------------------------------------------------------------------------
% Resumo
% -----------------------------------------------------------------------------

\begin{resumo}

    O resumo consiste em uma apresentação concisa de um texto, destacando seus elementos mais importantes. Trata-se de um recurso extremamente útil para o leitor, pois permite que ele compreenda rapidamente o conteúdo principal da obra.
    Nesse sentido, o resumo deve enfatizar o objetivo, o método, os resultados e as conclusões do trabalho. A ordem e a extensão desses elementos variam de acordo com o tipo de resumo e com a forma como cada item é abordado no texto original.
    Sua estrutura deve consistir em uma sequência de frases concisas, com caráter afirmativo e sem enumeração de tópicos, recomendando-se a utilização de um único parágrafo. As palavras-chave devem ser apresentadas logo abaixo do resumo, antecedidas da expressão “palavras-chave” e finalizadas por ponto.
    É importante evitar o uso de símbolos e contrações que não sejam de uso corrente, bem como fórmulas, equações, diagramas e elementos semelhantes que não sejam absolutamente necessários. Caso seu uso seja imprescindível, esses elementos devem ser definidos na primeira vez em que aparecerem.
    Quanto à extensão, os resumos devem seguir alguns parâmetros: para trabalhos acadêmicos, como teses, dissertações e relatórios técnico-científicos, recomenda-se entre 150 e 500 palavras; para artigos de periódicos, entre 100 e 250 palavras; e, para indicações breves, entre 50 e 100 palavras.
    Por fim, como parte integrante do trabalho, o resumo deve ser seguido pelas palavras representativas de seu conteúdo, isto é, palavras-chave, redigidos no mesmo idioma do texto.

    \par\vspace{\baselineskip}

    \textbf{Palavras-chave}: Palavra-chave 1; Palavra-chave 2; Palavra-chave 3; Palavra-chave 4.
\end{resumo}

% Em uma Tese de Doutorado o resumo deve conter, no máximo, 500 palavras.
% Em uma Dissertação de Mestrado o resumo deve conter, no máximo, 250 palavras.
% Em um Projeto de Qualificação o resumo deve conter, no máximo, 200 palavras.

% -----------------------------------------------------------------------------
% Escolha de 3 a 6 palavras ou termos que melhor representam seu trabalho.
% As palavras-chave são utilizadas para indexação. A letra inicial de cada
% palavra deve estar em maiúsculas. As palavras-chave são separadas por ponto.
% -----------------------------------------------------------------------------
